\documentclass[a4paper, 12pt]{article}

\usepackage{utils}

\renewcommand*{\today}{01 April 2026}

\begin{document}

\hotbox{Probabilités 1}{CM 8}{\today}

\section{Couples de variables aléatoires}

Dans ce chapitre, on considèra $(\Omega, \mathcal{T}, \P)$ un espace probabilisé.

\subsection{Cas discrete}

\begin{definition}
    Soit $\funcdef{}{\Omega}{\R}{\omega}{(X(\omega), Y(\omega))}$ une application où $X$ et $Y$ sont des variables aléatoires.

    La loi de $(X, Y)$ est appelée la \textbf{loi conjointe} du couple $(X, Y)$ et est définie par
    $$
    \forall (x, y) \in \R^2, \quad P((X, Y) = (x, y)) := P(X = x \text{ et } Y = y)
    $$
\end{definition}

\end{document}